\begin{abstract}
This study investigates interpretable unsupervised machine learning for identifying financial peer groups and changes in their membership among JSE Top 40 equities. The analysis extends a single-window study to three cross-sections: 30 June 2025, 31 December 2025 and 30 June 2026. Each cross-section contains 41 listed instruments representing 40 company groups. Eight financial variables were transformed using a June 2025 scaling reference and analysed through UMAP dimensionality reduction followed by hierarchical clustering with average linkage and Manhattan distance. Eight clusters were fitted independently at each date, and stability was assessed over five random seeds. The UMAP partitions achieved Silhouette Scores of 0.685, 0.638 and 0.721 in the reduced representations, compared with 0.151, 0.124 and 0.267 in the scaled financial features. Clusters were matched between dates using Jaccard similarity and a Hungarian assignment. Among 40 continuing instruments, 15 changed aligned group in the first period and 11 in the second. Of the 39 instruments present at all three dates, 19 retained their aligned group and 20 changed at least once. Cluster profiles, company-level feature changes, surrogate SHAP explanations and a separate Mahalanobis distance screen were used to interpret the results. The findings show how repeated clustering can support equity screening by identifying continuing and changing financial relationships. Retrospective data, representation sensitivity and incomplete external validation limit the conclusions that can be drawn about economic anomalies or investment performance.
\end{abstract}

Keywords: unsupervised learning; clustering; repeated cross-sections; dimensionality reduction; interpretable artificial intelligence; JSE Top 40; financial data mining.
